\section{Conclusión}

El desarrollo de este Trabajo Práctico nos permitió investigar y entender a fondo el poder de Programación Dinámica. Con el diseño de la ecuación de recurrencia y la matriz de memoization, se pudo resolver un escenario donde cada decisión presente desencadena un efecto dominó en el futuro, garantizando matemáticamente el máximo valor acumulado posible para Sophia en cada partida.

Al contrastar esta solución con el enfoque Greedy, mientras que solución avariciosa tenía un costo rapidísimo en $\mathcal{O}(n)$ usando memoria constante, falla al caer en ciertas ``trampas'' perdiendo la visión del problema completo frente a variabilidades grandes. En cambio, la Programación Dinámica exige pagar un precio alto: un sacrificio computacional cuadrático tanto en tiempo ($\mathcal{O}(n^2)$) como en espacio ($\mathcal{O}(n^2)$) para guardar la ``foto o recuerdos'' de la experiencia en una matriz. Pero nos da una garantía de optimalidad inquebrantable.

En conclusión, esta forma nos demuestra que la mejor solución global rara vez se encuentra eligiendo el camino más fácil y evidente en el presente. Resolver problemas complejos de verdad requieren la capacidad de sacrificar el beneficio inmediato, pensar en las ramificaciones del futuro, y tener la memoria suficiente para no recalcular los mismos errores del pasado.
